\section{Các tham số tự do của lời giải thích}
\label{sec:ch03-cac-tham-so-tu-do}

\index{tham số của LIME}%
Mục trước cho một phương trình, nhưng phương trình chưa tự xác định một lời
giải thích. Trước hết, người dùng phải quyết định biểu diễn $x'$ và phép dựng lại
$z$. Hai lựa chọn ấy quyết định thành phần nào được gọi tên và đầu vào nào
được đưa vào mô hình. Sau đó, phải chọn phân phối sinh nhiễu, hàm khoảng cách
$D$, độ rộng $\sigma$ của \tn{hàm nhân}{kernel}, số mẫu $N$ và độ dài $K$ của
lời giải thích. Bảy chỗ ấy là cách sách gom lại; bài báo không đếm chúng thành
một danh sách. Với phần lớn trong số đó, bài báo chốt sẵn một cách làm cho
riêng các thí nghiệm của mình, chẳng hạn lấy các thành phần khác không của $x'$
đều theo phân phối đều, dùng khoảng cách cosine cho văn bản và Euclid cho ảnh,
và giữ $K$ là một hằng số, để việc khảo sát các giá trị khác cho công trình sau.
Cái bài báo không làm là biện minh cho những lựa chọn ấy như thể chúng là điều
duy nhất làm được~\autocite{p09lime}.

Mỗi lựa chọn đứng thay cho một mệnh đề về hành vi của $f$ mà người dùng ngầm
chấp nhận khi đặt nó. Chọn $\sigma$ nhỏ là cho rằng $f$ chỉ gần tuyến tính
trong một bán kính rất hẹp quanh $x$, nên phép khớp ưu tiên các mẫu rất gần
$x$; chọn $\sigma$ lớn là cho rằng dạng tuyến tính còn đúng ở xa hơn, nên các
mẫu xa cũng góp trọng số đáng kể. Chọn $K$ nhỏ là cho rằng vài thành phần đủ
để mô tả hành vi của $f$ trong lân cận: lời giải thích dễ đọc hơn, nhưng có thể
bỏ mất một tương tác cần thiết. Tăng $N$ thì không mang mệnh đề nào như thế:
nó chỉ thay dữ liệu dùng để ước lượng $g$, chứ không làm cho $g$ tự nhiên trở
thành cơ chế thật của $f$.

Vì vậy, không nên gọi các lựa chọn này là \tn{siêu tham số}{hyperparameter}
thuần túy rồi để chúng ở ngoài phần diễn giải. Khi một báo cáo nói một từ, một super-pixel hay một
\tn{đặc trưng}{feature} quan trọng, người đọc cần biết lời giải thích đã cho phép những can
thiệp nào, đã coi mẫu nào là gần, và đã ép mô hình thay thế chỉ giữ bao nhiêu
thành phần. Trọng số của $g$ chỉ có nghĩa khi ba lựa chọn ấy được nêu ra cùng
với nó.

